\usepackage{geometry}
\usepackage{enumitem}
\usepackage{hyperref}

\usepackage{amsthm}
\usepackage{amsmath,mathtools}
\usepackage{amssymb}
\usepackage{bm}
\usepackage[per-mode=symbol]{siunitx}
\usepackage{tikz}
\usepackage{newpxmath}

\newcommand{\ee}{e} % 自然常数
\newcommand{\ii}{i} % 虚数单位
\newcommand{\dd}{\mathop{}\! d} % 微分算子
\newcommand{\K}{\vvmathbb K} % 数域
\newcommand{\N}{\vvmathbb N} % 自然数集
\newcommand{\Z}{\vvmathbb Z} % 整数集
\newcommand{\Q}{\vvmathbb Q} % 有理数集
\newcommand{\R}{\vvmathbb R} % 实数集
\newcommand{\C}{\vvmathbb C} % 复数集
\newcommand{\B}{\vvmathbb B} % 开球
\newcommand{\E}{\vvmathbb E} % 期望
\newcommand{\PP}{\vvmathbb P} % 概率
\newcommand{\power}{\mathcal P} % 幂集
\newcommand{\topo}{\mathcal T} % 拓扑
\newcommand{\nbr}{\mathcal N} % 邻域族
\newcommand{\onbr}{\mathcal O} % 开邻域族
\newcommand{\cnbr}{\mathcal F} % 闭邻域族
\newcommand{\basis}{\mathcal B} % 拓扑基
\newcommand{\subbasis}{\mathcal S} % 拓扑子基
\newcommand{\mat}{\bm} % 矩阵和向量
\newcommand{\pto}{\xlongrightarrow{p}} % 依概率收敛
\newcommand{\dto}{\xlongrightarrow{d}} % 依分布收敛
\newcommand{\coloniff}{:\iff} % 当且仅当的定义
\renewcommand{\implies}{\Longrightarrow} % 蕴含
\renewcommand{\impliedby}{\Longleftarrow} % 被蕴含
\renewcommand{\iff}{\Longleftrightarrow} % 等价
\DeclareMathOperator{\im}{im} % 像
\DeclareMathOperator{\ran}{ran} % 值域
\DeclareMathOperator{\id}{id} % 恒同映射
\DeclareMathOperator{\pr}{pr} % 投影
\DeclareMathOperator{\diam}{diam} % 集合的直径
\DeclareMathOperator{\vol}{vol} % 体积
\DeclareMathOperator{\supp}{supp} % 支撑集
\DeclareMathOperator{\nullity}{null} % 零化度
\DeclareMathOperator{\rank}{rk} % 秩
\DeclareMathOperator{\trace}{tr} % 迹
\DeclareMathOperator{\adj}{adj} % 伴随
\DeclareMathOperator{\spec}{spec} % 谱
\DeclareMathOperator{\var}{var} % 方差
\DeclareMathOperator{\cov}{cov} % 协方差
\let\oldRe\Re
\let\Re\relax
\DeclareMathOperator{\Re}{\oldRe} % 实部
\let\oldIm\Im
\let\Im\relax
\DeclareMathOperator{\Im}{\oldIm} % 虚部
\DeclareMathOperator{\Funct}{Funct} % 函数集合
\DeclareMathOperator{\Hom}{Hom} % 同态
\DeclareMathOperator{\End}{End} % 自同态
\DeclareMathOperator{\Iso}{Iso} % 同构
\DeclareMathOperator{\Aut}{Aut} % 自同构
\DeclareMathOperator{\Span}{span} % 张成
\DeclareMathOperator*{\upperlim}{\overline{\lim}} % 上极限
\DeclareMathOperator*{\underlim}{\underline{\lim}} % 下极限
\renewenvironment{cases}{
    \left\{\begin{array}{cl}
}{
    \end{array}\right.
}

\makeatletter

\ctexset{
    chapter={
        name={},
        number={第\arabic{chapter}章}
    },
    subsection={
        numbering=false
    }
}

\newtheoremstyle{my_def}
                {\thm@preskip}
                {\thm@postskip}
                {\slshape}
                {}
                {\bfseries}
                {\@addpunct{.}}
                {\thm@headsep}
                {}
\newtheoremstyle{my_rem}
                {\thm@preskip}
                {\thm@postskip}
                {}
                {}
                {\bfseries}
                {\@addpunct{.}}
                {\thm@headsep}
                {}

% \theoremstyle{plain}
\theoremstyle{my_def}
\newtheorem{definition}{定义}[section] % def
\newtheorem{axiom}[definition]{公理} % ax

\newtheorem{proposition}{命题}[section] % prop
\newtheorem{theorem}[proposition]{定理} % thm
\newtheorem{lemma}[proposition]{引理} % lem
\newtheorem{corollary}[proposition]{推论} % cor

% \theoremstyle{definition}
\theoremstyle{my_rem}
\newtheorem{remark}[proposition]{注} % rem
\newtheorem{example}[proposition]{例} % eg
\newtheorem{exercise}{}[section] % ex
% \renewcommand{\theexercise}{\arabic{exercise}}

\renewenvironment{proof}[1][\proofname]{\par
    \pushQED{\qed}
    \normalfont \topsep6\p@\@plus6\p@\relax
    \trivlist
    \item\relax
        {\bfseries
        #1\@addpunct{.}}\hspace\labelsep\ignorespaces
}{
    \popQED\endtrivlist\@endpefalse
}

\definecolor{winered}{RGB}{127.5, 0, 0}
\hypersetup{
    breaklinks,
    unicode,
    linktoc=page,
    bookmarksnumbered=true,
    bookmarksopen=true,
    colorlinks,
    % linkcolor=winered,
    plainpages=false,
    pdfstartview=Fit,
    pdfborder={0 0 0}
}

% Cover Styles ===============================
% 需要的宏包
\RequirePackage{geometry}
\RequirePackage{hyperref}
\RequirePackage{tikz}
\usetikzlibrary{calc}
% 颜色
\definecolor{coverlinecolor}{rgb}{255, 134, 24}
% 中文字体
\ifcsname heiti\endcsname
    \newcommand{\cbfseries}{\heiti}
\else
    \newcommand{\cbfseries}{\bfseries}
\fi
\ifcsname kaishu\endcsname
    \newcommand{\citshape}{\kaishu}
    \newcommand{\cnormal}{\kaishu}
\else
    \newcommand{\citshape}{\itshape}
    \newcommand{\cnormal}{\normalfont}
\fi
% 标签文字
\newcommand{\authorname}{\citshape 作者：}
\newcommand{\institutename}{\citshape 组织：}
\newcommand{\datename}{\citshape 时间：}
\newcommand{\versionname}{\citshape 版本：}
% 对外接口命令
\newcommand{\subtitle}[1]{\gdef\@subtitle{#1}}
\newcommand{\institute}[1]{\gdef\@institute{#1}}
\newcommand{\logo}[1]{\gdef\@logo{#1}}
\newcommand{\cover}[1]{\gdef\@cover{#1}}
\newcommand{\extrainfo}[1]{\gdef\@extrainfo{#1}}
\newcommand{\version}[1]{\gdef\@version{#1}}
\newcommand\bioinfo[2]{\gdef\@bioinfo{{\citshape #1}：#2}}
% 重写 maketitle
\renewcommand{\maketitle}{
    \hypersetup{pageanchor=false}
    % \pagenumbering{Alph}
    \begin{titlepage}
        \newgeometry{margin=0in}
        \parindent=0pt
        % 顶部封面图
        \ifcsname @cover\endcsname
            \includegraphics[width=\linewidth]{\@cover}
        \else
            \@empty
        \fi
        % 彩色横条
        % \setlength{\fboxsep}{0pt}
        % \colorbox{coverlinecolor}{\makebox[\linewidth][c]{\shortstack[c]{\vspace{0.5in}}}}
        % 标题
        \vfill
        \vskip 15ex % \vskip-2ex
        \hspace{2em}
        \parbox{0.8\textwidth}{
            \bfseries\Huge
            \ifcsname @title\endcsname \cnormal\@title \fi
            \par
        }
        % 副标题 + 信息
        \vfill
        % \vspace{-1.0cm}
        \hspace{2.5em}
        \begin{minipage}[c]{0.67\linewidth}
            \color{darkgray}
            {\bfseries\Large
            \ifcsname @subtitle\endcsname \cnormal\@subtitle\\[2ex] \fi}
            \color{gray}
            % \normalsize
            {\renewcommand{\arraystretch}{1.2}
            \begin{tabular}{l}
                % \ifcsname @author\endcsname \authorname \@author\\ \fi
                \ifx\@author\empty\else \authorname\cnormal\@author\\ \fi
                \ifcsname @institute\endcsname \institutename\cnormal\@institute\\ \fi
                % \ifcsname @date\endcsname \@date\\ \fi
                \ifx\@date\empty\else \datename\cnormal\@date\\ \fi
                \ifcsname @version\endcsname \versionname\cnormal\@version\\ \fi
                \ifcsname @bioinfo\endcsname \cnormal\@bioinfo\\ \fi
            \end{tabular}}
        \end{minipage}
        % 右下角 logo
        \begin{minipage}[c]{0.27\linewidth}
            \begin{tikzpicture}[remember picture,overlay]
                \node[opacity=0.8,
                    anchor=south east,
                    outer sep=0pt,
                    inner sep=0pt] at
                    ($(current page.south east)+(-0.8in, 1.5in)$)
                    {\ifcsname @logo\endcsname \includegraphics[width=4.2cm]{\@logo} \fi};
            \end{tikzpicture}
        \end{minipage}
        % 底部说明文字
        \vfill
        \begin{center}
            \parbox[t]{0.7\textwidth}{\centering\citshape
                \ifcsname @extrainfo\endcsname \@extrainfo \fi
            }
        \end{center}
        \vfill
    \end{titlepage}
    \restoregeometry
    \thispagestyle{empty}
}

\makeatother